\chapter{基于规则的知识获取}
\label{chap:rule-based-knowledge-acquisition}

一条程序能够顺利执行，是否就说明它产生了可信知识？设规则看到“任命”便输出“任职变动”。它会同时命中“海岬科技任命林舟为首席架构师”“海岬科技未任命林舟”和“关于任命林舟的传言已被否认”。字面匹配在三句话中都执行成功，只有第一句话却直接支持目标事件。规则能否运行、规则是否触发、触发后结论是否符合任务定义，是三个不同问题。

本章讨论\term{基于规则的知识获取}（rule-based knowledge acquisition）：把可检查的条件、已有知识或约束写成规则，再由规则为样本产生任务标签、候选范围或关系信息。重点是“候选知识凭什么产生”，而不是假定规则已经正确。专家编写、从数据归纳和由模型生成回答的是\emph{规则怎样形成}；模式匹配、知识关联和约束推理回答的是\emph{规则依据什么产生知识}。同一条规则可以由模型提出，却以知识库事实作为判断依据；也可以由专家编写，却只根据样本中的数值阈值触发。这两个维度不能混为一种分类。

为了让不同规则具有可比的执行接口，设样本为$x$，规则版本为$v$，适用前提为$A_v(x)$，触发条件为$C_v(x)$，输出构造为$O_v(x)$。本章采用三值逻辑$\{\mathsf{T},\mathsf{F},\mathsf{U}\}$分别表示条件成立、不成立和无法确定，先计算$D_v(x)=A_v(x)\land C_v(x)$，再把单条规则的结果写成
\begin{equation}
  r_v(x)=
  \begin{cases}
    O_v(x), & D_v(x)=\mathsf{T},\\
    \text{不适用}, & A_v(x)=\mathsf{F},\\
    \text{未触发}, & A_v(x)=\mathsf{T}\ \text{且}\ D_v(x)=\mathsf{F},\\
    \text{弃权}, & \text{其他情况}.
  \end{cases}
  \label{eq:rule-based-knowledge-acquisition-interface}
\end{equation}
这里的“不适用”是语义状态，表示对象不在规则声明的范围；“未触发”表示适用对象没有满足触发条件；“弃权”是来源动作，表示适用前提或触发条件因依据不足而无法确定。三者都不自动成为目标类别的负例。只有另一条规则明确检查了负类定义并输出负标签，才能把相应样本记作负例。

输出$O_v(x)$也不必是单个类别。它可以是“任职变动”这一候选标签，可以是$\{\text{管理岗位},\text{技术岗位}\}$这样的候选范围，也可以是“两个实体不得取同一身份”这样的关系约束。无论输出形式如何，规则都应同时留下命中的输入位置、使用的外部事实、经过的推理步骤或未能判断的原因。表~\ref{tab:rule-based-knowledge-acquisition-record}用同一条教学规则说明最低限度的记录内容。

\begin{table*}[htbp]
  \centering
  \begingroup
  \setlength{\tabcolsep}{4pt}
  \begin{tabularx}{\linewidth}{>{\raggedright\arraybackslash}p{19mm}>{\raggedright\arraybackslash}p{32mm}>{\raggedright\arraybackslash}p{49mm}>{\raggedright\arraybackslash}X}
    \toprule
    字段 & 作用 & 教学取值 & 缺少时无法判断的问题 \\
    \midrule
    输入对象 & 确定规则检查什么 & 一句中文新闻及其中的实体提及 & 规则究竟作用于整篇文档、句子还是实体对 \\
    适用前提 & 限定允许使用规则的范围 & 已完成中文切分、实体类型识别和否定范围分析 & 解析缺失时的结果是否仍在规则声明范围内 \\
    触发条件 & 说明何时产生输出 & “任命”连接组织与人员，且不在否定或被否认的引语中 & 命中由哪个条件造成，新增例外改变了什么 \\
    输出 & 说明产生何种候选知识 & 候选关系$(\text{林舟},\text{任职于},\text{海岬科技})$ & 返回的是类别、候选集合还是关系约束 \\
    无输出状态 & 区分范围、沉默与弃权 & 范围外记不适用，条件为假记未触发，条件未知时弃权 & 空输出表示范围外、未命中，还是依据不足 \\
    命中证据 & 保留从输入到输出的路径 & 关键词跨度、两个实体标识及局部依存路径 & 程序成功是否等于样本语义支持结论 \\
    适用范围及版本 & 支持更新与追溯 & 新闻语体；规则\texttt{appointment-v1}；知识库2026-09版 & 词表、解析器或外部知识更新后哪些旧结果需要重算 \\
    \bottomrule
  \end{tabularx}
  \endgroup
  \caption{一条知识获取规则需要记录的内容。教学规则可以返回关系候选，也可以按任务改为类别或候选集合；字段的作用是把语义状态、来源动作、执行状态、知识依据和版本范围分开。}
  \label{tab:rule-based-knowledge-acquisition-record}
\end{table*}

有了共享接口，下面可以按知识产生的主要依据比较三类方法。模式匹配检查样本是否满足可观察条件；知识关联检查样本能否对应已有实体、概念或事实；约束推理检查候选答案之间必须满足什么关系。三类方法可以出现在同一条规则中，但每一步使用了什么依据、可能在哪一步失败，仍应分别记录。

\section{基于模式匹配的规则利用}
\label{sec:rule-based-knowledge-acquisition-pattern}

\term{模式匹配}（pattern matching）直接检查样本中可观察的符号或数值条件，并在条件成立时产生输出。Snorkel中的标注函数提供了这类接口的一种实现：函数可以利用模式、启发式规则或外部资源为数据点给出标签，也可以在没有足够依据时弃权\scite{ratner2017}

模式规则的输入是样本及其可计算表示，输出是类别、候选范围或弃权，证据则是命中的跨度、数值与中间条件。失败可能发生在三个位置：输入没有所需字段，条件可算但不成立，条件成立却不能支持目标语义。前两种情况决定规则是否触发，第三种情况决定触发后的知识是否有依据。随着输入格式、领域用语和任务定义变化，模式本身及其前处理都需要维护。

\subsection{基于符号模式的规则利用}
\label{sec:rule-based-knowledge-acquisition-symbolic-pattern}

文本、日志和结构化记录中经常存在可直接观察的符号线索。关键词检查某个词是否出现，正则表达式进一步限制相邻字符与位置，句式模板或局部结构模式则要求若干成分以规定关系出现。它们的共同点是：规则依据当前样本中的符号结构触发，不需要先证明一个外部事实。

仍看“任命”规则。输入是已经切分的句子，匹配范围是句内谓词及其局部参数，触发条件要求出现任命谓词、施事是组织、受事是人员，并且谓词不受否定词支配，也不只出现在随后被否认的引语中。输出是任职变动候选及关系三元组，条件缺失或无法确定时弃权。对“海岬科技任命林舟为首席架构师”，执行过程依次记录“任命”的字符跨度、海岬科技的组织跨度、林舟的人员跨度以及三者的局部结构；这些位置共同组成命中证据。

只查关键词时，三段教学文本都会命中。加入“未”位于谓词否定范围内的条件，可以排除直接否定；但引用“任命林舟”的报道随后被否认时，局部窗口仍可能看见肯定形式。规则必须继续检查引语边界及报道者对引语的态度，或在无法可靠分析时弃权。图~\ref{fig:rule-based-knowledge-acquisition-pattern-rule}把这种收窄画成三层条件。新增条件改变的是规则适用范围，并没有仅凭条件更多就证明输出更准确。

\begin{figure*}[htbp]
  \centering
  \begin{tikzpicture}[
  x=1cm,
  y=1cm,
  font=\figurefont\footnotesize,
  panel/.style={draw=BookRule,line width=0.8pt,fill=white},
  text box/.style={draw=BookMuted,line width=0.8pt,fill=BookPaper,minimum height=10mm,text width=43mm,align=left,inner sep=2mm},
  condition/.style={draw=FigureModel,line width=1pt,fill=FigureModel!10,minimum height=9mm,text width=26mm,align=center,inner sep=1.5mm},
  state/.style={draw=BookMuted,line width=0.7pt,dashed,fill=white,minimum height=8mm,text width=16mm,align=center,inner sep=1mm},
  output/.style={draw=FigureOutput,line width=1pt,fill=FigureOutput!7,minimum height=9mm,text width=42mm,align=center,inner sep=1.5mm},
  trigger/.style={draw=FigureModel,line width=1pt,fill=FigureModel!10,minimum height=8mm,text width=18mm,align=center,inner sep=1mm},
  main/.style={-{Stealth[length=2mm]},draw=BookInk,line width=1.1pt},
  branch/.style={-{Stealth[length=1.8mm]},draw=BookMuted,line width=0.8pt,dashed},
  edge label/.style={fill=white,text=BookMuted,inner sep=1pt,font=\figurefont\scriptsize}
]
  \fill[white] (0,0) rectangle (16.6,6.8);
  \draw[panel] (0,0) rectangle (5.15,6.8);
  \draw[panel] (5.4,0) rectangle (11.05,6.8);
  \draw[panel] (11.3,0) rectangle (16.6,6.8);

  \node[anchor=north west,font=\figurefont\bfseries] at (0.2,6.6) {(a) 相同字面命中};
  \node[text box,anchor=north west] at (0.25,5.95) {肯定：海岬科技\textbf{任命}林舟。\\[-1pt]\textcolor{BookTeal}{字面命中}};
  \node[text box,anchor=north west] at (0.25,4.2) {否定：海岬科技\textbf{未任命}林舟。\\[-1pt]\textcolor{BookTeal}{字面命中}};
  \node[text box,anchor=north west] at (0.25,2.45) {引用后否认：“\textbf{任命}林舟”的传言被否认。\\[-1pt]\textcolor{BookTeal}{字面命中}};
  \node[anchor=south,text=BookMuted,font=\figurefont\scriptsize,align=center] at (2.575,0.25) {匹配成功只说明找到了跨度};

  \node[anchor=north west,font=\figurefont\bfseries] at (5.6,6.6) {(b) 组合条件与三值状态};
  \node[condition] (keyword) at (7.2,5.45) {出现目标谓词？};
  \node[condition] (roles) at (7.2,3.75) {组织与人员角色匹配？};
  \node[condition] (context) at (7.2,2.05) {不在否定或被否认引语中？};
  \node[state] (miss) at (9.8,5.45) {未触发\\条件为假};
  \node[state] (unknown) at (9.8,3.75) {弃权\\条件未知};
  \node[state] (excluded) at (9.8,2.05) {未触发\\例外成立};
  \node[trigger] (triggered) at (7.2,0.75) {规则触发};
  \draw[main] (keyword) -- node[edge label,left] {成立} (roles);
  \draw[main] (roles) -- node[edge label,left] {成立} (context);
  \draw[branch] (keyword.east) -- (miss.west);
  \draw[branch] (roles.east) -- (unknown.west);
  \draw[branch] (context.east) -- (excluded.west);
  \draw[main] (context) -- node[edge label,left] {成立} (triggered);

  \node[anchor=north west,font=\figurefont\bfseries] at (11.5,6.6) {(c) 任务所需输出};
  \node[output,anchor=north] (class) at (13.95,5.75) {类别候选\\“任职变动”};
  \node[output,anchor=north] (relation) at (13.95,3.95) {关系候选\\$(\text{林舟},\text{任职于},\text{海岬科技})$};
  \node[output,anchor=north,draw=BookAmber,fill=BookAmber!8] (abstain) at (13.95,2.15) {弃权\\条件为假或无法确定};
  \draw[main] (triggered.east) -- ++(3.85,0) |- (class.west);
  \draw[main] (triggered.east) -- ++(3.55,0) |- (relation.west);
  \node[anchor=south,text=BookAmber,font=\figurefont\bfseries\scriptsize,align=center] at (13.95,0.25) {未触发 $\neq$ 负例};
\end{tikzpicture}

  \caption{从字面命中到带例外的组合规则。左面板中相同谓词出现在肯定、否定和被否认的引语里；中面板的每个判断都显式保留$\mathsf{T}$、$\mathsf{F}$、$\mathsf{U}$三个出口，分别继续检查、记为未触发和弃权；右面板说明规则可以输出类别候选、候选范围或弃权。未触发和弃权都不是目标类别的负例。}
  \label{fig:rule-based-knowledge-acquisition-pattern-rule}
\end{figure*}

结构信息越丰富，前处理依赖也越多。字面关键词只依赖规范化与切分，却容易受多义词影响；带位置条件的正则表达式能约束局部顺序，却仍难处理跨句否定；句式模板可以检查实体角色和依存路径，但其结果又依赖实体识别、句法分析及领域句式。维护成本不能只按规则文本的长度计算，还要包括词形变化、拼写误差、分词方式和解析器版本带来的变化。

符号规则可以由专家直接写出，也可以从语料中的重复结构归纳，还可以让模型提出候选模板。这些方式改变的是候选规则从哪里来，不改变验证过程。每个候选都要回到原始样本，经历“检查支持样本和反例—定位误触发原因—收窄适用范围”的过程。若删去否定条件会新增哪些命中，加入实体类型又会放弃哪些原本可判断的样本，都应随版本记录。

执行检查之外，还要把“模式正确执行”与“命中后语义成立”分开。程序可能准确找到了每个“任命”跨度，却因为任务只标注已经生效的任职关系而把尚未生效的公告误作事实；程序也可能因解析失败而没有触发，但样本实际属于目标类。前一种是结论不符合任务定义，后一种是覆盖不足。如何用统一指标估计这类可靠性由第~\ref{chap:knowledge-quality-assessment-correction}章讨论，本节只要求证据路径可追溯。

\subsection{基于数值条件的规则利用}
\label{sec:rule-based-knowledge-acquisition-numeric-pattern}

有些判断材料不是词语，而是计数、比例、距离或测量值。\term{数值条件规则}（numeric-condition rule）先把输入转换为可比较的数值，再用阈值或区间产生候选判断。计算结果精确，只说明在给定输入与公式下比较没有算错；阈值为什么合理、适用于谁，仍需单独说明。

考虑一个明确标为假设的温度教学规则。任务约定：设备读数不低于$80\,^\circ\mathrm{C}$时输出“高温候选”，处于$[78,80)\,^\circ\mathrm{C}$时因测量误差影响而弃权，低于$78\,^\circ\mathrm{C}$时这条规则不触发。输入还必须带有单位、采集时间和传感器状态。于是$82\,^\circ\mathrm{C}$触发候选，$79\,^\circ\mathrm{C}$弃权，缺少单位同样弃权。$176\,^\circ\mathrm{F}$换算为$80\,^\circ\mathrm{C}$后恰好落在非严格下界上并触发；若只复制数值$176$而把单位误当成摄氏度，比较运算虽能完成，结论却来自错误输入。

一个完整的数值规则因此包含四步：计算特征，统一单位或基准，执行严格或非严格比较，再把比较状态映射为输出。多个区间必须说明端点归属和是否重叠；缺失值、溢出值和不可比较值必须有明确去向。不确定区间也不能任意添加，它应对应已知的测量误差、舍入范围或证据限制。若误差来源未知，可以保守弃权，却不能把随意选择的缓冲带解释为经过验证的置信区间。

阈值的形成方式决定应记录什么依据。按任务定义给出的阈值属于标签口径，任务定义改变时应整体重标；由领域知识确定的范围需要记录适用环境和标准版本；根据少量样本调出的经验阈值则依赖开发数据。经验阈值必须把调参数据与独立检查数据隔离，否则规则可能只记住开发样本。群体构成、采集设备或时间分布变化后，也要重新检查原阈值是否仍覆盖同一种对象。

固定阈值并不等于零维护成本。系统仍要计算特征、查询基准、换算单位、处理缺失值和监测输入变化。若基准随地区、设备或人群不同，规则还要先选择正确基准；选错基准与比较符号写错是两种故障。规则记录应让维护者能够区分阈值来源错误、特征计算错误和样本超出适用范围。

\subsection{基于组合条件的规则利用}
\label{sec:rule-based-knowledge-acquisition-composite-pattern}

真实规则通常要把多个基本条件组合起来。对条件$K$、$E$和$N$，可把“出现任命谓词”“实体角色匹配”“没有否定或被否认引语”分别记为三值量，再用$K\land E\land N$描述组合。只要有一项为假，合取条件为假；没有假值但至少一项未知时，结果为未知。这样，“实体识别器没有输出”不会被悄悄当成“实体类型不符合”。

组合方式还决定执行语义。\term{规则列表}（rule list）按优先级依次检查，首个触发项可以覆盖后续一般规则；\term{规则树}（rule tree）沿条件分支前进，每条路径都应到达输出或弃权；互不设优先级的\term{规则集合}（rule set）则分别执行各条规则，保留所有原始输出。列表中的默认分支不是天然负类，除非任务定义明确规定“前面所有条件不成立即可判负”。树中的条件缺失也不能随意走“否”分支，因为未知与假值可能导向不同处理。

短路执行可以节省计算，却不应删除证据。例如，关键词条件为假时无需继续运行昂贵的句法解析，但仍要记录规则因关键词未命中而停止；实体类型未知时则应记录未知来源，而不是伪装成关键词规则未触发。例外覆盖也要可见：若“否定语境”规则优先于一般任命规则，最终弃权记录中应同时保留一般模式命中和例外条件成立。

从数据归纳或由模型生成条件组合后，需要检查组合是否只记住开发样本。可以逐项删除条件，观察触发范围扩大到哪些新样本；也可以检查两个条件是否始终同时出现，从而暴露冗余。条件增加通常会缩小覆盖范围，却不必然提高可靠性，因为新增条件本身也可能错误。组合越深，调试每条路径、解释默认分支和维护前处理依赖的成本越高。

本节关注一条复合规则内部怎样执行。若多条独立规则分别输出不同标签，不能通过暗中调整执行顺序把冲突抹掉；应保留每条规则的版本、证据与原始输出，再交给第~\ref{chap:knowledge-fusion}章讨论的多来源融合过程。

\section{基于知识关联的规则利用}
\label{sec:rule-based-knowledge-acquisition-association}

样本中的字面条件有时不足以给出任务知识。例如，“林舟”这个名字本身不能说明指的是哪一位同名作者；“首席架构师”与“管理岗位”的关系也依赖采用哪套概念体系。\term{知识关联}（knowledge association）先把样本与已有实体、概念或事实对应起来，再利用对应结果产生候选知识。

这里必须检查两个层次：被查询的已有知识是否成立，以及这项知识是否适用于当前样本。实体库可能含有过时属性，概念体系可能采用不同分类目的，知识库关系即使真实也可能没有在当前句子中表达。表~\ref{tab:rule-based-knowledge-acquisition-association}比较三类关联的输入、依据和主要错误位置；它们可以在同一任务中共同使用，不是互斥算法。

\begin{table*}[htbp]
  \centering
  \begingroup
  \setlength{\tabcolsep}{2.5pt}
  \begin{tabularx}{\linewidth}{>{\raggedright\arraybackslash}p{17mm}>{\raggedright\arraybackslash}p{25mm}>{\raggedright\arraybackslash}p{26mm}>{\raggedright\arraybackslash}p{30mm}>{\raggedright\arraybackslash}p{29mm}>{\raggedright\arraybackslash}X}
    \toprule
    关联类型 & 样本中已知内容 & 查询的知识对象 & 需要成立的对应 & 产生的候选知识 & 主要歧义 \\
    \midrule
    实体匹配 & 名称、别名、标识及上下文 & 带唯一标识的实体及其属性 & 当前提及确实指向该实体，属性在当前时间和任务下适用 & 实体类别、属性或后续查询入口 & 同名不一定是同一实体；无候选与多候选不同 \\
    概念关联 & 词语、短语及局部用法 & 词典、本体或分类体系中的概念 & 表达与概念同一，且层级继承规则适用于目标属性 & 概念标签、上位类别或候选范围 & 相似词不一定是同一概念；上位概念通常不能确定唯一细类 \\
    关系匹配 & 已消歧实体对及证据范围 & 知识库中的有向关系事实 & 实体、方向、时间与关系定义一致，当前文本确实表达该关系 & 句子级关系候选或提及组级监督 & 共现不等于表达关系；知识缺失不等于关系为假 \\
    \bottomrule
  \end{tabularx}
  \endgroup
  \caption{三类知识关联的输入、依据与错误位置。表中每一行都要求同时保存样本侧证据和知识侧版本，避免只留下最终标签而无法判断错误发生在哪一步。}
  \label{tab:rule-based-knowledge-acquisition-association}
\end{table*}

\subsection{基于实体匹配的规则利用}
\label{sec:rule-based-knowledge-acquisition-entity-matching}

本章沿用\term{实体匹配}这一上位名称，表示在样本对象与已有实体之间建立对应；对于文本，本节具体采用\term{实体链接}（entity linking）接口：把样本中的一个提及对应到知识库中带唯一标识的对象。它解决“这段文字指的是谁或什么”，而不是仅判断两个字符串是否相同。输入包括提及跨度、名称、别名、标识符和上下文，输出可以是唯一实体、多个候选、没有候选或因证据不足而弃权。

设教学实体库中有两位都叫“林舟”的作者：实体\texttt{E-017}与海岬科技相关，实体\texttt{E-204}与东岭大学相关。句子“林舟介绍了海岬科技的新系统”会先按名称生成两个候选，再用同句组织、作者主页标识或文档来源消歧。名称完全相同并没有直接给出唯一实体；若上下文只写“林舟发表了报告”，合理输出是$\{\texttt{E-017},\texttt{E-204}\}$而不是任意选择一个。若实体库中根本没有这位作者，则是“无候选”；若有多个候选却无法区分，则是“多候选”。两个状态需要不同的补充信息。

精确标识符匹配通常证据最强，但前提是标识确实属于当前提及；别名表能处理简称和旧名，却需要维护名称生效时间；模糊字符串匹配可以容忍拼写变化，也会扩大误配范围。模型辅助消歧可以利用更广上下文，但它仍只负责选择关联对象。实体链接研究通常把候选生成与基于上下文的候选选择分开，这一区分有助于定位“没找到实体”和“选错实体”两类故障\scite{knowledge-acquisition-milne2008}

链接到正确实体以后，属性转移仍需检查任务定义。例如，实体库把\texttt{E-017}的职业记录为“软件工程师”，当前任务却标注文档中当时担任的职务；静态职业属性不能直接覆盖特定时间的任职事件。事实本身错误、提及链接错误和属性不适用于当前标注对象，应分别记录。规则至少要保存实体库版本、属性生效时间、候选列表、消歧证据和最终选择。

实体词表与别名规则可以由专家维护，也可以从历史匹配数据归纳，或由模型建议新增别名。这些仍是规则形成方式。无论别名来自哪里，都需要检查它是否跨越了不同实体、是否只在特定时间有效，以及移除该别名会影响哪些旧结果。热门实体拥有更多别名和上下文记录，可能因此更容易被匹配；这种覆盖差异也应进入后续可靠性检查。

\subsection{基于概念关联的规则利用}
\label{sec:rule-based-knowledge-acquisition-concept-association}

实体回答“具体是哪一个对象”，概念回答“这个表达属于哪类含义”。\term{概念关联}（concept association）把样本中的词语或短语映射到词典、本体或分类体系中的概念标识，再按明确的上下位关系获得类别或候选范围。词语形式只是入口，真正用于推导的是概念标识及体系中的关系。

设一套教学岗位体系把“首席架构师”映射到概念\texttt{role:\allowbreak{}technical-\allowbreak{}lead}，该概念的上位类是\texttt{role:\allowbreak{}management-\allowbreak{}position}。文本确认细类以后，可以按体系中的“属于”边产生上位类别“管理岗位”。反向推理却不成立：只知道某人处于管理岗位，不能推出其唯一细类是首席架构师，因为同一上位类还包含其他岗位。若任务只需要粗类别，上位概念是可用输出；若任务需要细类，就应保留下位候选集合。

同义词、缩写、近义概念和上下位概念需要不同规则。同义词可以指向同一概念标识；缩写还要检查领域和上下文；近义词只表示含义接近，不能自动合并；上下位词则保留层级方向。WordNet将同义词组织为同义词集，并用语义关系连接这些概念，说明词面集合与概念关系可以分层记录\scite{knowledge-acquisition-miller1995}

概念映射表、本体关系和推导规则应分开维护。映射表回答“文本表达指向哪个概念”，本体回答“概念之间有什么关系”，推导规则回答“哪些关系允许转成任务标签”。属性也不能沿层级任意复制。例如，“属于管理岗位”可以由细类向上继承，但某个具体岗位的薪酬范围、任期或资格要求未必适用于全部父子概念。

不同知识体系可能按职责、组织级别或行业标准分类，同一个词在两套体系中因分类目的不同而落在不同位置。规则必须记录采用的体系与版本。概念合并、拆分或父子关系调整后，应沿映射和推导记录定位受影响的旧结果。缺少概念关联只说明当前体系没有提供映射，不能据此推出样本不属于目标类别；跨来源概念对齐的通用问题留给第~\ref{chap:knowledge-fusion}章。

\subsection{基于关系匹配的规则利用}
\label{sec:rule-based-knowledge-acquisition-relation-matching}

已有知识还可以直接记录实体之间的事实。\term{关系匹配}（relation matching）把样本中的实体对与知识库关系对齐，再产生关系标签或成组监督。关系通常有方向和时间，例如$(\text{林舟},\text{任职于},\text{海岬科技})$不同于反向三元组，也不能自动推广到事实生效前后的全部文本。

\term{远程监督}（distant supervision）是关系匹配的代表性用法。Mintz等人的设定把知识库实体关系与文本中共同出现的实体对对齐，从而在没有逐句人工关系标签时构造训练实例\scite{knowledge-acquisition-mintz2009} 这种对齐提供的是启发式监督，而不是逐句事实证明。

设知识库中的构造事实为“林舟任职于海岬科技”。句子甲“海岬科技任命林舟为首席架构师”确实表达了与任务关系相符的任职事件；句子乙“林舟在论坛上批评了海岬科技的发布流程”只同时提到两者。若把知识库事实无条件复制给所有共同出现的句子，句子乙就会得到没有文本证据的正标签。图~\ref{fig:rule-based-knowledge-acquisition-relation-path}在实体消歧、关系定义对应和证据范围判断之后，分别保留句子级候选与实体对提及组级候选。

\begin{figure*}[htbp]
  \centering
  \begin{tikzpicture}[
  x=1cm,
  y=1cm,
  font=\figurefont\footnotesize,
  source/.style={draw=FigureInput,line width=0.9pt,fill=FigureInput!9,minimum height=10mm,text width=45mm,align=left,inner sep=2mm},
  fact/.style={draw=FigureModel,line width=1pt,fill=FigureModel!10,minimum height=10mm,text width=53mm,align=center,inner sep=2mm},
  process/.style={draw=BookMuted,line width=0.8pt,fill=BookPaper,minimum height=10mm,text width=24mm,align=center,inner sep=1.5mm},
  decision/.style={draw=BookAmber,line width=1pt,fill=BookAmber!9,minimum height=12mm,text width=25mm,align=center,inner sep=1.5mm},
  result/.style={draw=FigureOutput,line width=1pt,fill=FigureOutput!7,minimum height=11mm,text width=28mm,align=center,inner sep=1.5mm},
  flow/.style={-{Stealth[length=2mm]},draw=BookInk,line width=1.1pt,rounded corners=1mm},
  evidence/.style={-{Stealth[length=1.8mm]},draw=FigureModel,line width=0.9pt,dashed,rounded corners=1mm},
  edge label/.style={fill=white,text=BookMuted,inner sep=1pt,font=\figurefont\scriptsize}
]
  \fill[white] (0,0) rectangle (16.7,6.8);
  \node[fact] (kb) at (8.3,6.1) {已有构造事实\\$(\text{林舟/E-017},\text{任职于},\text{海岬科技/E-003})$};

  \node[source,anchor=east] (sent-a) at (4.9,4.65) {句甲：海岬科技任命林舟为首席架构师。};
  \node[source,anchor=east] (sent-b) at (4.9,2.75) {句乙：林舟在论坛上批评了海岬科技的发布流程。};

  \node[process] (entity) at (6.35,3.7) {实体消歧\\确认同一实体对};
  \node[process] (mapping) at (9.25,3.7) {关系定义对应\\方向与时间};
  \node[decision] (scope) at (12.2,3.7) {当前句是否\\表达该事实？};

  \node[result] (positive) at (15.2,4.65) {句子级候选\\句甲：任职于};
  \node[result,draw=BookMuted,dashed,fill=white] (cooccur) at (15.2,2.75) {不产生句子正例\\句乙：仅共同出现};
  \node[result,text width=42mm] (bag) at (12.2,1.05) {提及组级候选\\至少一条提及可能表达关系};

  \draw[flow] (sent-a.east) -- ++(0.35,0) |- (entity.west);
  \draw[flow] (sent-b.east) -- ++(0.35,0) |- (entity.west);
  \draw[flow] (entity) -- (mapping);
  \draw[flow] (mapping) -- (scope);
  \draw[flow] (scope.east) -- ++(0.2,0) |- (positive.west);
  \draw[flow] (scope.east) -- ++(0.2,0) |- (cooccur.west);
  \draw[flow] (positive.south) -- ++(0,-0.35) -| (bag.east);

  \draw[evidence] (kb.south) -- (mapping.north);
  \draw[evidence] (kb.south east) -- ++(2.9,-0.55) |- (scope.north);

  \draw[flow] (0.35,0.65) -- (1.25,0.65);
  \node[anchor=west,text=BookMuted,font=\figurefont\scriptsize] at (1.35,0.65) {处理步骤};
  \draw[evidence] (2.95,0.65) -- (3.85,0.65);
  \node[anchor=west,text=BookMuted,font=\figurefont\scriptsize] at (3.95,0.65) {事实的证据关联};
\end{tikzpicture}

  \caption{从已有事实到文本候选关系的关联路径。实体和关系均为构造例子；实线箭头表示处理步骤，虚线表示知识库事实提供的证据关联。句甲和句乙沿独立路径完成实体消歧、关系定义对应与证据范围判断：句甲表达目标关系，句乙仅共同提及两个实体。两条提及路径与知识库事实共同进入提及组级候选，因此组内可以保留“至少一条提及可能有证据”的判断，却不能把每句都改成确定正例。}
  \label{fig:rule-based-knowledge-acquisition-relation-path}
\end{figure*}

完整对齐过程先确认实体及方向，再检查知识库关系与任务关系的定义是否一致，随后确定证据范围是当前句、整段还是同一实体对的多条提及。Riedel等人把“实体对存在关系”和“某个提及是否表达该关系”作为不同变量，并采用至少一次表达的假设，正是为了放松“每个共同提及都表达关系”的过强假设\scite{knowledge-acquisition-riedel2010}

因此，同一实体对的多条提及适合表示成一个组。组级关系候选可以表示“这些提及中至少一条可能支持关系”，却不等于每句都是正例。后续模型若要在组内识别具体证据，需要另行建立提及级推断机制；本章只要求规则输出注明粒度，并保留哪些句子、实体标识和知识库事实参与了关联。

错误也应沿路径分解。实体消歧错误会把无关对象连在一起；关系事实过时会在错误时间提供监督；文本没有表达关系则是证据范围不成立。知识库没有记录某个关系时，还存在开放世界问题：缺失可能来自覆盖不足，不能自动作为负例。规则应输出未知、候选关系或由明确排除证据支持的负标签，而不是把“未查到”与“已证明不存在”混为一谈。

知识覆盖、关系类别不平衡、热门实体的高出现频率和重复转载都会改变生成候选的分布。它们不一定由规则执行日志直接暴露，所以至少要保存知识库版本、事实来源、实体对、提及列表和去重标识，以便后续融合与可靠性检查追查来源。

\section{基于约束推理的规则利用}
\label{sec:rule-based-knowledge-acquisition-constraint}

模式匹配从样本条件出发，知识关联从已有对象和事实出发；还有一类规则不直接指定每个位置的答案，而是规定答案之间必须满足什么关系。\term{约束推理}（constraint-based inference）把未知标签放入候选空间，用逻辑、结构或领域条件排除不可能的组合，必要时再按软偏好选择候选。

设全部候选答案为$\mathcal{Y}$，硬约束为$c_1,\ldots,c_m$，则可行集合为
\begin{equation}
  \mathcal{F}
  =
  \{\bm{y}\in\mathcal{Y}:c_k(\bm{y})=\mathsf{T},\ k=1,\ldots,m\}.
  \label{eq:rule-based-knowledge-acquisition-feasible-set}
\end{equation}
推理可能得到一个解、多个解或空集。一个解表示约束恰好确定了候选；多个解表示依据只能缩小范围；空集表示已知事实和约束相互冲突，或输入有误。三种都是需要显式报告的求解状态。尤其不能在多解时任取一个当作知识，也不能在无解时悄悄放松约束。

\subsection{基于逻辑约束的规则利用}
\label{sec:rule-based-knowledge-acquisition-logical-constraint}

\term{逻辑约束}（logical constraint）用蕴含、互斥、等价等关系描述命题之间允许怎样共同取值。硬约束完全排除违反它的候选；\term{软约束}（soft constraint）允许违反，但为违反分配代价。软约束只有与权重或其他偏好结合时，才能在多个逻辑可行答案之间作选择。

设二值命题$p$表示“属于技术管理岗位”，$q$表示“属于管理岗位”，教学体系给出硬约束$p\Rightarrow q$。四种赋值中，$(p,q)=(1,0)$被排除，其他三种仍可行。若再人为设置软约束代价$L(\bm{y})$，可以选出代价最小的一个候选，但“代价最小”只表示符合当前偏好，不等于已经证明它是真实答案。图~\ref{fig:rule-based-knowledge-acquisition-feasible-set}用可复算的教学数值显示这一区别。

\begin{figure*}[htbp]
  \centering
  \begin{tikzpicture}[
  x=1cm,
  y=1cm,
  font=\figurefont\footnotesize,
  candidate/.style={draw=BookMuted,line width=0.8pt,fill=BookPaper,minimum width=30mm,minimum height=13mm,align=center,inner sep=1.5mm},
  selected/.style={candidate,draw=FigureModel,line width=1.2pt,fill=FigureModel!10},
  premise/.style={draw=FigureInput,line width=0.9pt,fill=FigureInput!9,minimum width=28mm,minimum height=10mm,align=center,inner sep=1.5mm},
  constraint/.style={draw=BookAmber,line width=1pt,fill=BookAmber!9,minimum width=31mm,minimum height=10mm,align=center,inner sep=1.5mm},
  combine/.style={circle,draw=BookInk,line width=0.9pt,fill=white,minimum size=7mm,inner sep=0pt},
  empty/.style={draw=BookAmber,line width=1.2pt,fill=BookAmber!9,minimum width=32mm,minimum height=11mm,align=center},
  flow/.style={knowledge arrow,draw=BookInk,line width=1.1pt},
  soft/.style={text=BookMuted,font=\figurefont\scriptsize}
]
  \fill[white] (0,0) rectangle (16.6,6.5);
  \node[anchor=north west,font=\figurefont\bfseries] at (0,6.25) {(a) 硬约束筛选与软偏好};
  \node[constraint,minimum width=44mm] (hard) at (8.3,5.75) {硬约束：$p\Rightarrow q$};

  \node[candidate] (c00) at (1.8,4.25) {$(p,q)=(0,0)$\\{\scriptsize 可行；$L=1.0$}};
  \node[candidate] (c01) at (5.65,4.25) {$(p,q)=(0,1)$\\{\scriptsize 可行；$L=0.6$}};
  \node[candidate,draw=BookAmber,fill=BookAmber!7] (c10) at (9.5,4.25) {$(p,q)=(1,0)$\\{\scriptsize 违反硬约束}};
  \node[selected] (c11) at (13.35,4.25) {$(p,q)=(1,1)$\\{\scriptsize 可行；$L=0.2$（最小）}};
  \draw[BookAmber,line width=1.2pt] ([xshift=-10mm,yshift=-4mm]c10.center) -- ([xshift=10mm,yshift=4mm]c10.center);
  \draw[BookAmber,line width=1.2pt] ([xshift=-10mm,yshift=4mm]c10.center) -- ([xshift=10mm,yshift=-4mm]c10.center);
  \node[soft,anchor=north] at (8.3,3.42) {硬约束留下三个答案；软代价只负责排序};

  \draw[BookRule,line width=0.8pt] (0,2.95) -- (16.6,2.95);
  \node[anchor=north west,font=\figurefont\bfseries] at (0,2.75) {(b) 新增互相矛盾的已知条件};
  \node[premise] (p-one) at (2.2,1.65) {观测：$p=1$};
  \node[premise] (q-zero) at (7.2,1.65) {观测：$q=0$};
  \node[constraint] (imply) at (12.2,1.65) {硬约束：$p\Rightarrow q$};
  \node[combine] (and) at (7.2,0.35) {$\land$};
  \node[empty] (none) at (12.4,0.35) {无可行解\\报告冲突};
  \draw[flow] (p-one.south) -- ++(0,-0.35) -| (and.west);
  \draw[flow] (q-zero.south) -- (and.north);
  \draw[flow] (imply.south) -- ++(0,-0.35) -| (and.east);
  \draw[flow] (and) -- (none);
\end{tikzpicture}

  \caption{约束缩小可行答案并产生冲突的教学示意。上面板从两个二值命题的四种赋值出发，硬约束$p\Rightarrow q$排除$(1,0)$，软约束代价再在三个可行候选中给出唯一最小值；这些代价是人为设置，不是统计估计。下面板同时加入观测$p=1$、观测$q=0$和同一蕴含后，可行集合变为空集。逻辑可行、偏好最小和真实正确是三种不同判断。}
  \label{fig:rule-based-knowledge-acquisition-feasible-set}
\end{figure*}

一般地，令第$j$条软约束的违反代价为$\ell_j(\bm{y})\geq0$，其中零表示不违反，正值表示违反程度；令对应权重为$w_j\geq0$，可以在硬约束可行集上求
\begin{equation}
  \widehat{\bm{y}}
  \in
  \argmin_{\bm{y}\in\mathcal{F}}
  \sum_j w_j\ell_j(\bm{y}).
  \label{eq:rule-based-knowledge-acquisition-soft-constraint}
\end{equation}
这个表达式把两种结论分开：$\bm{y}\in\mathcal{F}$说明没有违反硬约束，目标值较小说明按当前权重更受偏好。约束推断可以通过整数或线性规划整合局部预测与任务条件，Roth和Yih给出了自然语言任务中的代表性形式\scite{knowledge-acquisition-roth2004} 将逻辑公式作为可违反的加权关系时，权重的统计含义和落地规模还需要另外建模；Markov逻辑网络展示了这类软逻辑的一种概率化实现\scite{pgm-energy-models-richardson2006}

规则形成时最容易发生的错误之一是反转蕴含。由“首席架构师属于技术负责人”只能写出$\text{首席架构师}\Rightarrow\text{技术负责人}$，不能反推所有技术负责人都是首席架构师。另一类错误是省略适用对象，使只针对正式岗位的互斥条件误用于临时角色。专家、数据归纳或模型生成的候选约束，都要核对方向、量词、对象和例外。

推理记录应保留使用了哪些前提、排除了哪些候选以及哪些软约束贡献了代价。约束更新以后，维护者才能定位受影响结果。求解规模、软权重设定和冲突解释都是成本；若约束主要用于综合多个现成来源，而不是从已知前提产生新候选，则问题转入第~\ref{chap:knowledge-fusion}章的融合范围。

\subsection{基于结构约束的规则利用}
\label{sec:rule-based-knowledge-acquisition-structural-constraint}

标签往往不是彼此独立的。\term{结构约束}（structural constraint）利用序列、树、图或成组记录的组织关系，限制合法标签组合。结构可以来自任务格式、标注规范或输入数据，也可能由实体匹配等前置步骤构造。满足结构只说明组合合法，通常不足以证明每个标签语义正确。

以序列标记为例，本节采用BIO约定：\texttt{B-X}表示类型X片段的开头，\texttt{I-X}表示延续同类型片段，\texttt{O}表示片段外；\texttt{I-X}只能跟在\texttt{B-X}或\texttt{I-X}之后。早期文本分块工作已经把片段结构编码成逐词标签，说明局部标签需要服从边界表示约定\scite{knowledge-acquisition-ramshaw1995}

对“林舟 加入 海岬 科技”，每个位置可以先由局部规则产生候选集合。若逐点选择得到\texttt{B-PER, I-ORG, B-ORG, I-ORG}，第二个位置的\texttt{I-ORG}前面不是组织片段开头，因此整体非法。图~\ref{fig:rule-based-knowledge-acquisition-sequence-constraint}中的合法行\texttt{B-PER, O, B-ORG, I-ORG}满足转移条件，但它只是一种合法候选；例如把“海岬科技”全部标为\texttt{O}也满足BIO形式，仍可能不符合真实语义。

\begin{figure*}[htbp]
  \centering
  \begin{tikzpicture}[
  x=1cm,
  y=1cm,
  font=\figurefont\footnotesize,
  token/.style={draw=FigureInput,line width=0.9pt,fill=FigureInput!9,minimum width=30mm,minimum height=8mm,align=center},
  tag/.style={draw=FigureModel,line width=0.9pt,fill=FigureModel!9,minimum width=30mm,minimum height=8mm,align=center},
  bad tag/.style={tag,draw=BookAmber,fill=BookAmber!9},
  choices/.style={draw=BookMuted,line width=0.7pt,dashed,fill=white,minimum width=34mm,minimum height=9mm,align=center,font=\figurefont\scriptsize},
  allowed/.style={knowledge arrow,draw=BookInk,line width=1pt},
  invalid/.style={knowledge arrow,draw=BookAmber,line width=1pt,dashed},
  edge label/.style={fill=white,text=BookMuted,inner sep=1pt,font=\figurefont\scriptsize}
]
  \fill[white] (0,0) rectangle (16.85,6.3);
  \node[anchor=east,font=\figurefont\bfseries] at (1.25,5.85) {词元};
  \foreach \x/\word in {3.1/林舟,7.3/加入,11.5/海岬,15.1/科技} {
    \node[token] at (\x,5.85) {\word};
  }

  \node[anchor=east,font=\figurefont\bfseries,text=BookAmber] at (1.25,4.45) {局部};
  \node[tag] (bad1) at (3.1,4.45) {\texttt{B-PER}};
  \node[bad tag] (bad2) at (7.3,4.45) {\texttt{I-ORG}};
  \node[tag] (bad3) at (11.5,4.45) {\texttt{B-ORG}};
  \node[tag] (bad4) at (15.1,4.45) {\texttt{I-ORG}};
  \draw[invalid] (bad1) -- (bad2);
  \draw[allowed] (bad2) -- (bad3);
  \draw[allowed] (bad3) -- (bad4);
  \node[circle,draw=BookAmber,fill=white,text=BookAmber,line width=0.8pt,minimum size=4mm,inner sep=0pt,font=\figurefont\bfseries\scriptsize] at (5.2,4.45) {$\times$};

  \node[anchor=east,font=\figurefont\bfseries,text=FigureModel] at (1.25,3.0) {合法};
  \node[tag] (good1) at (3.1,3.0) {\texttt{B-PER}};
  \node[tag] (good2) at (7.3,3.0) {\texttt{O}};
  \node[tag] (good3) at (11.5,3.0) {\texttt{B-ORG}};
  \node[tag] (good4) at (15.1,3.0) {\texttt{I-ORG}};
  \draw[allowed] (good1) -- node[edge label,above] {允许} (good2);
  \draw[allowed] (good2) -- node[edge label,above] {允许} (good3);
  \draw[allowed] (good3) -- node[edge label,above] {允许} (good4);

  \node[anchor=east,font=\figurefont\bfseries,text=BookMuted] at (1.25,1.45) {候选};
  \node[choices] at (3.1,1.45) {\{\texttt{B-PER}, \texttt{O}\}};
  \node[choices] at (7.3,1.45) {\{\texttt{O}, \texttt{I-ORG}\}};
  \node[choices] at (11.5,1.45) {\{\texttt{B-ORG}, \texttt{O}\}};
  \node[choices] at (15.1,1.45) {\{\texttt{I-ORG}, \texttt{O}\}};
  \node[anchor=south,text=BookMuted,font=\figurefont\scriptsize,align=center] at (8.4,0.2) {结构约束排除非法组合，但合法序列仍可能不止一个};
\end{tikzpicture}

  \caption{局部标签与整体结构的约束关系。采用BIO边界约定：\texttt{B-X}开始类型X片段，\texttt{I-X}只能延续同类型的\texttt{B-X}或\texttt{I-X}，\texttt{O}位于片段外。上行的\texttt{B-PER}\,$\to$\,\texttt{I-ORG}违反转移条件；下行是一种合法候选。每个位置下方仍保留局部候选集合，说明联合约束会排除组合，却不保证只剩一个合法序列，也不把下行称为真实答案。}
  \label{fig:rule-based-knowledge-acquisition-sequence-constraint}
\end{figure*}

局部传播可以在相邻位置间反复删除不可能标签；当一个选择会同时影响多个位置、树分支或图节点时，则需要动态规划、搜索或约束求解共同决定。边界包含、父子标签和跨记录身份一致性也属于结构约束，但它们检查的关系不同。尤其要区分稳定身份与随时间变化的属性：同一实体跨记录应保持实体标识一致，不代表其岗位、位置或状态永远不变。

结构本身也需要证据。文档层级可能由文件格式直接给出，跨记录的“同一对象”可能来自实体匹配，图的邻接边则可能由距离或关系规则构造。缺边会阻断本应发生的传播，错边会把无关对象绑在一起，错误边界会让合法标签也显得冲突。推理不能弥补输入结构缺少依据；结构来源和版本必须与标签证据一起保存。

结构大小与连接复杂度决定求解负担。长序列的局部转移可以用动态规划处理，稠密图上的全局条件却可能需要更大的搜索空间。采用近似方法时，应明确哪些冲突尚未解决、输出是完整解还是局部一致结果。无论求解多精确，结构合法都不等于语义真实。

\subsection{基于领域约束的规则利用}
\label{sec:rule-based-knowledge-acquisition-domain-constraint}

有些候选要受到守恒关系、数量平衡、范围限制或业务一致性的约束。\term{领域约束}（domain constraint）把这些关系写成未知标签或数值必须满足的条件，再由已知观测缩小候选。它与一般逻辑约束的区别不在求解器，而在约束依据来自特定领域，因而必须记录变量含义、适用环境、单位和容差。

考虑一个完全构造的物料平衡例子。初始量$s_0=100\,\mathrm{kg}$，流入量$q_{\mathrm{in}}=30\,\mathrm{kg}$，结束量$s_1=108\,\mathrm{kg}$，待求流出量为$q_{\mathrm{out}}$。在“没有其他流入、流出或损失”的教学前提下，精确关系
\begin{equation}
  s_1=s_0+q_{\mathrm{in}}-q_{\mathrm{out}}
  \label{eq:rule-based-knowledge-acquisition-balance}
\end{equation}
给出$q_{\mathrm{out}}=22\,\mathrm{kg}$。这个唯一值来自精确等式和完整观测共同作用，并不是任何物料系统都天然满足的无条件结论。

若测量误差只允许检查
\[
  \left|s_1-(s_0+q_{\mathrm{in}}-q_{\mathrm{out}})\right|
  \leq 2\,\mathrm{kg},
\]
则规则只能推出$q_{\mathrm{out}}\in[20,24]\,\mathrm{kg}$。若流入量缺失，候选范围还会更宽；若观测为$s_1=140\,\mathrm{kg}$且仍要求$q_{\mathrm{out}}\geq0$，上述“无其他来源”前提会导致无解。唯一解、候选范围和观测冲突必须作为不同状态返回，并记录单位、容差和求解精度。

规则来源决定结论强度。按任务定义成立的业务恒等式可以作为硬约束；领域规律若只在封闭、稳态或理想条件下成立，就必须把这些条件写入适用前提；从数据中观察到的经验关系通常更适合作为软约束，并保留估计误差与更新条件。表~\ref{tab:rule-based-knowledge-acquisition-domain-constraint}用同一教学例子比较三种表达。

\begin{table*}[htbp]
  \centering
  \begingroup
  \setlength{\tabcolsep}{2.5pt}
  \begin{tabularx}{\linewidth}{>{\raggedright\arraybackslash}p{23mm}>{\raggedright\arraybackslash}p{32mm}>{\raggedright\arraybackslash}p{25mm}>{\raggedright\arraybackslash}p{26mm}>{\raggedright\arraybackslash}p{32mm}>{\raggedright\arraybackslash}X}
    \toprule
    依据类型 & 前提 & 可观察量 & 误差或容差 & 允许推出的结论 & 前提失效时的处理 \\
    \midrule
    按任务定义成立的关系 & 任务把各量定义为完整且互斥的记账项 & 假设$s_0=100$、$q_{\mathrm{in}}=30$、$s_1=108$，单位均为$\mathrm{kg}$ & 定义中无额外容差 & 由精确等式得到$q_{\mathrm{out}}=22\,\mathrm{kg}$ & 报告输入违反任务定义或无解，不暗中增加遗漏项 \\
    理想条件下的领域规律 & 假设系统封闭，期间没有未记录损失，传感器误差不超过$2\,\mathrm{kg}$ & 同上，并核验封闭状态和传感器版本 & $\pm2\,\mathrm{kg}$ & 得到$q_{\mathrm{out}}\in[20,24]\,\mathrm{kg}$的可行范围 & 条件未知时弃权；条件不成立时改用包含损失项的模型 \\
    数据归纳的经验关系 & 关系只在指定设备、时间和样本范围内观察到 & 同类测量及形成规则的数据版本 & 用软代价及已核验的残差范围表示 & 只提供偏好排序或异常候选，不宣称守恒定律 & 分布或设备变化后重新估计；证据不足时保留多个候选 \\
    \bottomrule
  \end{tabularx}
  \endgroup
  \caption{领域约束的前提决定结论强度。表中数值均为单位一致的假设教学数据；精确等式、带容差范围和经验软关系分别支持唯一值、可行区间和偏好候选，不能相互替代。}
  \label{tab:rule-based-knowledge-acquisition-domain-constraint}
\end{table*}

专家提出领域约束时，需要核对变量是否与数据字段同义；模型提出候选关系时，还要检查它是否把相关性误作恒等式。环境改变以后，原本成立的封闭条件、误差范围或业务流程可能失效。开发成本不仅包括写出公式，还包括对齐观测、换算单位、建模误差、解释无解和维护前提版本。

领域约束可以产生候选，却不能独立证明输入测量可靠。第~\ref{chap:knowledge-quality-assessment-correction}章将统一讨论所得知识的质量、错误检测和修正；本节只要求约束强度与其依据相符，并在前提不满足时返回范围、冲突或弃权，而不是强行输出单值。

\section{本章小结}
\label{sec:rule-based-knowledge-acquisition-summary}

规则的形成方式与知识产生依据是两个正交问题。专家编写、数据归纳和模型生成说明候选规则怎样得到；模式匹配、知识关联和约束推理则分别以可观察条件、已有实体或事实、候选间必须满足的关系作为知识依据。同一条规则可以交叉使用多种机制，但每一步的输入、输出、证据与版本仍应分开记录。

模式匹配能够直接输出任务标签或候选关系，代价是持续维护符号变化、数值基准和组合分支。知识关联把样本连接到实体、概念与事实，能够复用已有知识，也引入消歧、版本、时间和证据范围问题。约束推理可以排除不合法答案、缩小候选集合或发现冲突，却只有在前提和结构有依据时才有意义；软约束选出的最小代价答案更不能自动等同于真实答案。

把本章的教学例子连起来，可以得到一条完整路径：符号模式在句子中找到“任命”及其实体角色，实体链接确认“林舟”和“海岬科技”的标识，关系事实提供任职候选，时间与结构约束再检查方向和有效期。对象不在规则范围时记为不适用；适用前提或触发条件未知时，规则应弃权；模式没有触发时，只能说当前规则未触发；约束留下多个解时，应保留候选范围；约束无解时，应报告冲突。只有明确检查了负类定义的规则，才能输出负例。

规则的明确形式使错误可以追溯到输入、前提、匹配、关联或求解步骤，却不保证结论适用于所有样本。值得维护的规则应具有可核验依据、清楚适用范围和可接受维护成本；覆盖不足时需要补充模式或知识，来源冲突时进入第~\ref{chap:knowledge-fusion}章，结论可靠性与修正则由第~\ref{chap:knowledge-quality-assessment-correction}章统一处理。
